% Esquisse d'un programme (cote 119, copy 2): archive cote 119, site page 113 (PDF page 114), typescript page 35.
% typescript. Draft compared: esquisse/pages/esq_119_037.tex (copy 1, site 37).
% AI-transcribed by Claude Opus 5.5 (claude-opus-5-5, adaptive thinking, effort high), Message Batches API, batch esq119_full (msgbatch_015xqDHkTumDysRF1qngZhRJ), 2026-11-13; not yet checked by eye.
\scanpage{119/113}
\mspage{35}

géométriques courantes, par exemple pour l'espace modulaire com-%
pact $\hat{M}_{g,\nu}$ avec sa ``stratification à l'infini'' canonique, donnée
par le diviseur à croisements normaux de Mumford-Deligne. Une
autre difficulté, moins sérieuse sans doute, c'est que le soi-%
disant ``espace'' modulaire est en fait une \uline{multiplicité} - techni-%
quement, cela s'exprime surtout par la nécessité de remplacer
l'ensemble d'indices $I$ pour les strates par une \uline{catégorie} (es-%
sentiellement finie) d'indices, en l'occurence celle des
``graphes MD'', qui ``paramètrent'' les ``structures combinatoires''
possibles d'une courbe stable de type $(g,\nu)$. Ceci dit, je puis
affirmer que la théorie de dévissage général, spécialement déve-%
loppé sous la pression du besoin de \uline{cette} cause, s'est révélé
en effet un guide précieux, conduisant à une compréhension pro-%
gressive, d'une cohérence sans \unsure{failles,de|failles, de} certains aspects essen-%
tiels de la tour de Teichmüller (c'est à dire, essentiellement
de la``structure à l'infini'' des groupes de Teichmüller ordinaires).
C'est cette approche qui m'a conduit finalement, dans les mois
suivants, vers le principe d'une construction purement combina-%
toire de la tour des groupoïdes de Teichmüller, dans l'esprit es-%
quissé plus haut (cf. par. 2).

Un autre test de cohérence satisfaisant provient du
point de vue ``topossique''. En effet, mon intérêt pour les multi-%
plicités modulaires provenant avant tout de leur sens algébrico-%
géométrique et arithmétique, c'est aux multiplicités modulaires
\uline{algébriques}, sur le corps de base absolu $\mathbb{Q}$, que je me suis inté-%
ressé prioritairement, et à un ``dévissage'' à l'infini de leurs
groupes fondamentaux géométriques (i.e. des groupes de Teichmüller
\uline{profinis}) qui soit compatible avec les opérations naturelles de
$\mathrm{I}\!\Gamma = \mathrm{Gal}(\bar{\mathbb{Q}}/\mathbb{Q})$. Cela semblait exclure d'emblée la possibilité de me
référer à une hypothétique théorie de dévissage de structures
stratifiées dans un contexte de ``topologie modérée'' (ou même de
topologie ordinaire, cahin-caha), si ce n'est comme fil conducteur
entièrement heuristique. Dès lors se posait la question de tra-%
duire, dans le contexte des topos (en l'occurence les topos étales)

\margin{113}
